\begin{lemma}\label{lemma.bitLengthIsDegree}
 For $b \in \N$, \codehl{b.bit\_length()} $= \deg \binaryRepresentation(b) + 1$.
\end{lemma}
\begin{proof}
 For $b > 0$, \codehl{bit\_length} returns one more than the position of the highest set bit,
 which is $\deg \binaryRepresentation(b)$ by definition.
 For $b = 0$ it returns $0 = \deg 0 + 1$, by the convention $\deg 0 = -1$.
\end{proof}

So \codehl{b.bit\_length() - 1} is the degree of $b$, the position of its highest set bit.
The lowest set bit takes one step more.
For $b > 0$ write $\trailingZeros(b)$ for the number of zeros below the lowest set bit of $b$.

\begin{prop}\label{prop.lowestSetBit}
 For $b \in \N$ with $b > 0$, $b \bitAnd (-b) = 2^{\trailingZeros(b)}$.
 Hence $\trailingZeros(b)$ is \codehl{(b \& -b).bit\_length() - 1}.
\end{prop}
\begin{proof}
 Let $n = \trailingZeros(b)$, so that bit $k$ of $b$ is $0$ for $k < n$ and $1$ for $k = n$.
 Subtracting $1$ from $b$ flips exactly its trailing zeros and its lowest one:
 bit $k$ of $b - 1$ is $1$ for $k < n$, $0$ for $k = n$, and bit $k$ of $b$ for $k > n$.
 By Definition \ref{def.twosComplementSequence},
 $\twosComplement(-b) = \neg\,\binaryRepresentation(b-1)$,
 so bit $k$ of $-b$ is $0$ for $k < n$, $1$ for $k = n$,
 and the complement of bit $k$ of $b$ for $k > n$.
 Taking the minimum position by position gives $0$ below $n$, $1$ at $n$,
 and $f_k \wedge (1 - f_k) = 0$ above it:
 $b \bitAnd (-b)$ has exactly one set bit, at position $n$.
 The second claim is Lemma \ref{lemma.bitLengthIsDegree} applied to $2^{n}$.
\end{proof}

For $b = 360$, with the thirteen lowest bits written out and position $0$ on the right,
\begin{equation*}
 \begin{array}{rl}
  b & 0000101101000 \\
  \neg b & 1111010010111 \\
  -b = \neg b + 1 & 1111010011000 \\
  b \bitAnd (-b) & 0000000001000
 \end{array}
\end{equation*}
and $\trailingZeros(360) = 3$.
The two reverse steps clear a bit.
$b \oplus 2^{\trailingZeros(b)}$ clears the lowest set bit, since $x \oplus x = 0$;
and $b \oplus 2^{\deg b}$ clears the highest,
which is the step \codehl{bits \^{}= 1 << index} of the walk of
Listing \ref{listing.tickArrayWalk}.
Listing \ref{listing.lowestSetBit} checks the proposition on a thousand integers of two hundred
bits, against a count of the trailing zeros one bit at a time.

\begin{lstlisting}[caption={$b \bitAnd (-b) = 2^{\trailingZeros(b)}$.}, label=listing.lowestSetBit]
import random


def degree(b):
    return b.bit_length() - 1


def trailing_zeros(b):
    count = 0
    while b & 1 == 0:
        count += 1
        b >>= 1
    return count


for _ in range(1000):
    b = random.randint(1, 2**200)
    assert b & -b == 2 ** trailing_zeros(b)
    assert degree(b & -b) == trailing_zeros(b)
    assert (b ^ (b & -b)).bit_count() == b.bit_count() - 1
    assert (b ^ 1 << degree(b)).bit_length() < b.bit_length()
\end{lstlisting}

The case $b = 0$ is excluded, and it matters.
Then $b \bitAnd (-b) = 0$ and \codehl{(0).bit\_length()} is $0$,
so both formulae return $-1$:
in a book, a position below the edge of the band,
and so a plausible price returned without an error.
An empty side has to be tested for before either formula is applied,
which is the test \codehl{if not bits} opening the lookup of Listing \ref{listing.tickArrayWrite}.
